% chapters/01_terminology_and_map.tex
\chapter{术语与地图：推断期学习的概念版图}
\label{ch:terminology}

\section{本章想回答什么问题}

在深度学习文献中，``Train''、``Inference''、``Adaptation''、``ICL''、
``Test-Time Update'' 这些词经常混用。
更严重的是，``Test-Time Training'' 这个词在 2020 年和 2024 年指代的是两种
\emph{根本不同}的机制——它们的更新对象、更新方式、所需基础设施都截然不同。

本章的目标：在阅读任何后续内容之前，建立清晰的术语边界，防止混淆。

\section{最小例子}

考虑一个在 ImageNet 上训练好的分类模型 $f_\theta$，将其部署到一台自动驾驶汽车上。
\begin{itemize}
  \item \textbf{训练（Training）}：在大规模 ImageNet 数据上，通过反向传播更新 $\theta$。
        完成后，$\theta$ 被固定在服务器上。
  \item \textbf{推断（Inference）}：汽车遇到一个路牌图片，调用 $f_\theta(\text{路牌图片})$，
        得到预测结果。$\theta$ 不变。
  \item \textbf{测试时更新（Test-Time Update）}：如果系统在遇到路牌时
        先做 5 步自监督微调再预测——这就是"测试时更新"的通俗定义。
  \item \textbf{上下文学习（ICL）}：如果换成大语言模型，
        在提示词中给出一些路牌翻译样例，让模型直接推断——
        参数 $\theta$ \emph{完全不变}，但输出受到上下文影响。
\end{itemize}

\section{直觉解释：两条 TTT 线索的本质差异}

\subsection*{第一条线索：TTT 2020 风格（外科手术式微调）}

\citet{sun2020ttt} 提出的"Test-Time Training" 指的是：
当测试样本 $x_\text{test}$ 到来时，利用一个\textbf{自监督辅助任务}
对网络参数 $\theta$ 执行真实的梯度下降更新，然后再做主任务预测。

\begin{center}
\texttt{输入测试样本} $\to$ \texttt{自监督微调（修改 $\theta$）} $\to$ \texttt{主任务预测}
\end{center}

关键点：这需要调用标准的反向传播（\texttt{.backward()}）。
每处理一个测试样本，都要运行一次或多次梯度步骤。
这种方法\textbf{真实地改变了模型的持久化参数}。

\subsection*{第二条线索：TTT-Layer 2024 风格（内嵌计算式更新）}

\citet{sun2024ttt} 以及相关工作（见第 \ref{ch:ttt_layers} 章）
把"TTT" 重新定义为：
将序列模型的\textbf{隐藏状态} $s_t$ 本身看作一个微型学习器的参数，
在前向传播过程中，每个新 token 到来时就执行一次基于状态的梯度步骤。

\begin{center}
\texttt{前一隐藏状态 $s_{t-1}$} $\to$ \texttt{输入 $x_t$} $\to$
\texttt{内部梯度步骤（更新 $s_t$）} $\to$ \texttt{输出 $y_t$}
\end{center}

关键点：\emph{整个过程都在前向传播内完成}。
没有外部的 \texttt{.backward()} 调用。
模型的\textbf{持久化参数 $\theta$ 不变}；变化的是每次推断的隐状态。

\section{数学形式化：三类"可更新对象"}

我们用一个统一记号来区分不同层级的更新对象：

\begin{definition}[推断期更新的三个层级]
设模型为 $f$，持久化参数为 $\theta$，隐藏状态为 $s_t$，输入为 $x_t$。
\begin{enumerate}
  \item \textbf{第0层（不变）}：$y_t = f_\theta(x_t)$，$\theta$ 和 $s$ 均不变。
  \item \textbf{第1层（状态更新）}：$s_t = g_\theta(s_{t-1}, x_t)$，$y_t = h_\theta(s_t, x_t)$。
        $\theta$ 不变，$s_t$ 在序列内随时间演化。这是标准 RNN 的语言。
  \item \textbf{第2层（参数更新）}：$\theta_t = \theta_{t-1} - \eta \nabla_\theta \mathcal{L}(x_t; \theta_{t-1})$。
        持久化参数 $\theta$ 在测试阶段被更新。这是 TTT 2020 的语言。
\end{enumerate}
\end{definition}

\begin{remark}
TTT-Layer（第 \ref{ch:ttt_layers} 章）处于第2层和第1层之间的一个有趣位置：
它把\emph{隐状态}本身当作需要被"梯度下降"更新的内部参数，
但这个更新完全内嵌在前向传播中，不需要外部梯度流。
\end{remark}

\section{概念全图}

\begin{figure}[h]
\centering
% 图示占位符——见 figures/ 目录
\fbox{\parbox{0.9\textwidth}{\centering
\textbf{全书概念演化图（待插图）}\\[1ex]
TTT 2020 $\to$ MAML $\to$ Fast Weights $\to$ Linear Attention \\
$\to$ ICL as GD $\to$ TTT Layers $\to$ Titans $\to$ Nested Learning
}}
\caption{推断期学习的演化主线。箭头表示概念依赖或历史演化方向。}
\label{fig:concept_map}
\end{figure}

\section{和前后章节的关系}

本章是\textbf{导航章}，不引入新的数学内容。
下一章（第 \ref{ch:why_test_time} 章）从最直接的动机出发：
面对分布偏移，静态模型为什么会失效，以及 TTT 2020 如何响应这个挑战。

\begin{misconceptionbox}
\textbf{常见误解 1}：``Test-Time Training'' 就是指在测试时跑反向传播。

\textbf{纠正}：从 2024 年起，TTT 越来越多地指代\emph{隐状态级别的序列内更新}，
与传统意义上的"反向传播微调权重"是不同的机制。
始终注意论文是在使用哪种定义。
\end{misconceptionbox}

\begin{misconceptionbox}
\textbf{常见误解 2}：ICL（上下文学习）一定涉及参数更新。

\textbf{纠正}：ICL 的参数 $\theta$ 完全不变。
模型利用上下文 token 的注意力交互来"适应"——
这种适应是否等价于梯度下降是第 \ref{ch:icl} 章的核心问题，
但无论如何，这与 TTT 意义上的参数更新不同。
\end{misconceptionbox}

\section{练习题}

\begin{exercise}
区分以下操作属于第0层、第1层还是第2层更新：
(a) 标准 Transformer 做一次前向传播；
(b) LSTM 按时间步更新隐藏状态；
(c) MAML 的内循环梯度步骤；
(d) GPT-4 在提示中读取 5-shot 样例后回答问题。
\end{exercise}

\begin{exercise}
TTT 2020 需要在测试时调用 \texttt{backward()}，而 TTT-Layer 不需要。
从硬件（GPU 显存）的角度，这两种方案有什么本质不同？
哪种方案更适合对延迟敏感的实时推断？
\end{exercise}

\begin{exercise}
自查：在你读过的一篇论文或博客中，作者是否混用了"训练"和"测试时更新"这两个概念？
如果有，请用本章的框架将其重新表述。
\end{exercise}

\section{本章小结}

\begin{itemize}
  \item 推断期学习有三个层级：参数不变（标准推断）、隐状态更新（RNN/Attention）、参数更新（TTT）。
  \item "TTT" 这个词在 2020 年和 2024 年含义不同：一个是外部梯度微调权重，一个是前向传播内嵌的隐状态学习。
  \item ICL 不改变参数，但仍可被分析为某种形式的隐式适应——第 \ref{ch:icl} 章详述。
  \item 本书的演化主线是：从 TTT 2020 的直观动机出发，逐步揭示更深的机制，最终抵达 Nested Learning 的统一视角。
\end{itemize}
